\section{Physical water management: when does a neighboring hydrophobic material protect a catalyst?}
\label{sec:water}

\textbf{Proposed contribution.} Determine whether a small polydivinylbenzene (PDVB) addition protects an already conditioned cobalt catalyst against a subsequent water disturbance, and establish whether measured water and catalytic responses predict useful output under a new disturbance history. The outcome would be a controlled test at a defined catalyst state and disturbance, with a conditional predictive rule within the calibrated range. A unique microscopic explanation of active-site drying is a conditional extension. No experiments proposed here have been performed.

\subsection{Prior work, importance, and the remaining question}

Physical hydrophobic promotion is established. Fang and colleagues demonstrated PDVB promotion of Co--Mn carbide for syngas conversion, including mixing-distance and water-cofeed effects; after removing PDVB from the used mixture, the recovered carbide performed comparably to the original catalyst. Li and colleagues extended physical promotion to copper-catalyzed methanol synthesis.\cite{fang2022,li2023} These precedents rule out claiming physical mixing, post-use additive removal, cross-reaction transfer, or a contact-angle correlation as the original contribution.

The immediate platform is Fang and colleagues' Co/SiO\textsubscript{2} study of cobalt Fischer--Tropsch synthesis, which converts CO and hydrogen (syngas) to hydrocarbons while producing water. C\textsubscript{5+} denotes hydrocarbons containing at least five carbon atoms. With 0.025\,g PDVB per 0.5\,g catalyst, separately prepared 40--60 mesh granules maintained approximately 44\% CO conversion and 88--90\% C\textsubscript{5+} selectivity over 1200\,h at 220\,$^\circ$C and 2\,MPa. The quartz-diluted reference initially converted more CO but declined strongly; nominal bed volumes were matched. Their water-cofeed, heat-transfer, material-state, and mixing controls support a substantial functional effect.\cite{fang2026} The mechanistic interpretation nevertheless requires separating changes in retained water, water transport, reaction rate, and catalyst history.

For perspective, let $X_{\mathrm{CO}}$ be the fraction of inlet CO converted and $S_{\mathrm{C5+}}$ the published carbon-based selectivity, counting carbon atoms rather than product molecules. Integrating their published series over the common 25--585\,h interval gives
\begin{equation}
 Q^*=\int X_{\mathrm{CO}}(t)S_{\mathrm{C5+}}(t)\,dt
 =116.0\ \text{h (reference)},\quad 219.8\ \text{h (PDVB)}.
 \label{eq:water-output}
\end{equation}
This secondary calculation uses Fang's Source Data workbook, sheet ``Figure 1,'' with equal nominal CO feed. Conversion and selectivity are separately interpolated piecewise linearly, then their product is integrated on the combined time grid over 25--585\,h. The output-proxy ratio is 1.89, or 1.80 after charging the additional 5\% polymer mass; the latter excludes quartz mass. It is not a complete carbon balance, an uncertainty interval, or a lifetime extrapolation. Early operation before 25\,h is excluded. The practical question is therefore credible, but superiority to an optimized cobalt formulation or regeneration policy has not been established.\cite{fang2026}

Water history and pore environment also have direct precedents. Hanssen and colleagues already compared water pretreatment with subsequent cofeeding and a dry return on alumina-supported cobalt.\cite{hanssen1997} Wolf and colleagues used externally supplied water at very low conversion followed by a dry return to connect cobalt oxidation, sintering, and lost function on graphite. Paterson and colleagues showed that activation-water exposure affects subsequent cobalt performance.\cite{wolf2019,paterson2024} Yang and colleagues related silica-shell geometry to deactivation and distinguished water-related and heavy-hydrocarbon-related behavior.\cite{yang2020} More recently, Sengupta and colleagues reported that strongly retained oxygenates change catalyst water affinity during operation.\cite{sengupta2025} Thus a dry fresh-material contact angle cannot be assumed to represent a working catalyst. The proposed distinction is \emph{late addition after a common conditioning protocol}, followed by a controlled functional comparison and a prediction outside the calibration history.

\subsection{What a faster water transient can establish}

At equilibrium, surfaces can hold different amounts of water at the same chemical potential. Hydrophobicity alone cannot maintain a chemical-potential gradient across an otherwise equilibrated connected gas. Reaction, imposed flow, or another source or sink can maintain a gradient; wettability can then affect storage, liquid configuration, and the resistance to transport. This distinction preserves the reported promotion while rejecting equilibrium loading differences as a sufficient explanation of continuous drying.

A minimal reversible-storage balance makes the ambiguity explicit:
\begin{equation}
 B\frac{dc}{dt}=q-G(c-c_b),\qquad
 \tau=\frac{B}{G},\qquad c_{\mathrm{ss}}-c_b=\frac{q}{G}.
 \label{eq:water-pool}
\end{equation}
Here $c$ and $c_b$ are local and boundary concentrations in a common gas-referenced coordinate, $B$ is differential storage capacity (volume), $G$ is exchange conductance (volume/time), and $q$ is water production (amount/time). The expressions assume a locally linear, time-invariant compartment with a specified source. Decreasing $B$ shortens both filling and emptying without lowering the stationary source-generated excess. It can increase exposure earlier in an upward step while shortening the subsequent tail. Consequently, less storage need not improve durability. Increasing $G$ also shortens the transient but lowers the steady excess. The reported syngas-to-Ar decay times of approximately 80 and 175\,s cannot distinguish these alternatives: that switch also changes production and adsorbed species.\cite{fang2026}

Even absolute inventory plus several outlet transients need not identify the producing region. For two pools with concentrations $c_i$, capacities $B_i$, and conductances $G_i=\lambda B_i$, $\lambda$ is their common exchange rate constant (inverse time). Let $B_T=B_1+B_2$ and $W=B_1c_1+B_2c_2$. With total water production $q_T$, inlet concentration $c_{\mathrm{in}}$, a well-mixed gas volume $V$ at concentration $c_g$, and actual volumetric flow $F$,
\begin{align}
 \dot W&=\lambda(B_Tc_g-W)+q_T,\\
 V\dot c_g&=F(c_{\mathrm{in}}-c_g)+\lambda(W-B_Tc_g).
 \label{eq:water-two-pool}
\end{align}
Only total capacity and total source enter the measured outlet $c_g$. If all production is in pool 1, its steady pool-to-gas excess is $c_{1,\mathrm{ss}}-c_{g,\mathrm{ss}}=q_T/(\lambda B_1)$. Changing $B_1$ from $B_T/4$ to $3B_T/4$ reduces this excess threefold without changing any outlet response at identical aggregate initial conditions. Since $c_{g,\mathrm{ss}}-c_{\mathrm{in}}=q_T/F$, the total pool-to-inlet excess includes an unchanged $q_T/F$ term and does not generally scale threefold. This constructed example is not a model fitted to PDVB; it shows why a source-specific or spatial constraint is needed for an active-site claim. Isotope signals do not automatically supply that constraint.

The program therefore separates three levels of evidence: measured formulation response, prospective prediction of catalytic function, and identification of a local mechanism. A correct outlet fit meets only the first. A stationary transport explanation must also constrain the location of the source and relevant resistance. For example, changing an external resistance cannot remove an unchanged internal resistance. Changing granule size alone does not resolve their relative contributions.

\subsection{Experiment 1: qualify the platform and the intervention}

Reproduce the low-dose, separate-granule contrast with measured catalyst inventory, bed volume, temperature, pressure drop, carbon balance, and useful-product rate. Keep the source's distinct configurations separate: its imaging of a powder co-granulate and some high-dose controls do not define the low-dose longevity packing. The published restart procedure is incomplete: synthesis includes reduction and passivation, whereas Supplementary Note 10 calls the standard test hydrogen-pre-reduced without fully specifying that restart. Establish and document one reproducible activation sequence before interpreting the new comparison; do not silently import the 300\,$^\circ$C analytical pretreatment.\cite{fang2026}

The late intervention requires cooling and repacking. Use independently preweighed catalyst charges from the same preparation, each with identical initial quartz dilution that remains with the charge throughout handling. Condition without PDVB under syngas, initially using the platform's H\textsubscript{2}/CO/Ar ratio of 64/32/4 at 220\,$^\circ$C and 2\,MPa. The parent pilot selects a fixed duration beyond startup, using rate and selectivity drift over a defined observation window with tolerances set from pilot variability. Freeze the duration and complete feed, flow, temperature, and pressure schedule before randomization; report pre-handling rate/selectivity drift and retained-liquid observations. This defines an operationally conditioned age, not microscopic stationarity or cessation of changes in water affinity.

Randomize charges within preparation and run-order blocks to the four conditions in Table~\ref{tab:water-design}. Common protocol controls intended history; independent repetitions quantify the remaining state variability. Recover each complete catalyst-and-quartz zone under inert gas, preserve retained liquid, and use the original dry catalyst inventory as the mass denominator after demonstrating quantitative solid recovery. Record liquid loss separately. All charges undergo identical cooling, transfer, and restart; no solvent washing or new reduction is introduced between conditioning and challenge.

\begin{table}[ht]
\centering
\caption{The decisive late-additive comparison. Each condition has independent catalyst charges; repeated chromatograms are not independent replicates.}
\label{tab:water-design}
\begin{tabularx}{\linewidth}{@{}lXX@{}}
\toprule
After common conditioning & Water challenge & Time-matched control\\
\midrule
0.05\,g PDVB/g catalyst & Low--high--low water & Low water throughout\\
Matched quartz diluent & Low--high--low water & Low water throughout\\
\bottomrule
\end{tabularx}
\end{table}

Repack each whole recovered charge with either 0.05\,g PDVB/g catalyst or quartz occupying the same measured bulk volume, to equal final bed volumes; no initial diluent is selectively removed. Keep catalyst granule size fixed and measure actual void volume and pressure drop rather than assuming equal hydrodynamics. Include the initial and final packing, heat and flow behavior, and solid/liquid recovery in the handling qualification. This pilot compares restarted parent charges with an uninterrupted parent control. A reproducible shift can define the post-handling starting state; uncontrolled solid/liquid loss or erratic restart closes the late-addition route. Handling cannot recreate uninterrupted industrial operation, and this limitation remains part of the interpretation.

\subsection{Experiment 2: distinguish protection from an activity penalty}

Use one imposed-water challenge selected in a parent-only feasibility study. A starting design is 220\,$^\circ$C, 20\,bar total pressure, inlet $p_{\mathrm{H2}}=10$\,bar and $p_{\mathrm{CO}}=5$\,bar, with water raised from 0.5 to 4\,bar by replacing Ar at fixed inlet molar flow. These are design choices, not a demonstrated damage threshold for Co/SiO\textsubscript{2}. Select the plateau and recovery durations before the four-condition comparison: the plateau must outlast measured equilibration and leave measurable but noncatastrophic residual loss in the parent. If no suitable disturbance exists within the qualified operating range, the proposed damage test has failed qualification.

Target at most 2\% CO conversion in the more active arm. Assuming one water per CO converted, the generated/supplied molar-water ratio is
\begin{equation}
 \frac{\dot n_{\mathrm{water,gen}}}{\dot n_{\mathrm{water,in}}}
 \simeq\frac{X_{\mathrm{CO}}y_{\mathrm{CO,in}}}{y_{\mathrm{water,in}}}
 =0.025\ \text{at 4 bar water},\qquad 0.20\ \text{at 0.5 bar water}.
 \label{eq:water-source}
\end{equation}
Thus external water dominates the high plateau much more strongly than recovery. Determine actual ratios from balances, including CO\textsubscript{2}, oxygenates, gas contraction, and unresolved oxygen storage. Low conversion suppresses whole-bed reactant changes; it does not prove small intragranular gradients. Before fixing the flow, verify that a modest further flow increase does not materially change mass-normalized diagnostic rate. An order-of-magnitude planning calculation, retaining the source's nominal promoted CO consumption, requires about 68\,L\,(g catalyst h)$^{-1}$ to reach 2\% conversion at inlet CO fraction 0.25. For a 50\,mg charge this corresponds to about 3.4 standard L/h of gas and 0.55\,g/h supplied water at 20\% steam (0\,$^\circ$C volume convention). These are equipment-sizing estimates, not rate predictions. Reducing inventory can make the test practical without reducing production per individual granule.

For each charge measure pre-challenge and post-recovery C\textsubscript{5+} carbon-production rates, $r_{\mathrm{pre}}$ and $r_{\mathrm{post}}$, after defined settling windows. Let $R=r_{\mathrm{post}}/r_{\mathrm{pre}}$. The primary protection contrast is the interaction
\begin{equation}
 D=(R_{\mathrm{PDVB,challenge}}-R_{\mathrm{PDVB,sham}})
  -(R_{\mathrm{quartz,challenge}}-R_{\mathrm{quartz,sham}}).
 \label{eq:water-interaction}
\end{equation}
A positive $D$ means that PDVB reduces the fractional functional penalty associated with the challenge beyond time-matched drift. Report both pre-challenge and post-recovery absolute rates for all four conditions alongside $D$; a favorable fraction can conceal an inferior catalyst bed. The second decision is the direct difference in integrated C\textsubscript{5+} carbon output per complete bed volume over the same predefined restart--challenge--recovery horizon. That comparison charges the initial activity penalty and includes startup product inventories and collection delays. Report catalyst-plus-additive mass productivity separately. The first endpoint tests protection, and the second tests whether protection is useful on that horizon.

Choose the number of independent blocks from pilot variability and a predeclared minimum consequential difference in retention and output; no fixed percentage is justified before that information exists. Report intervals incorporating preparation and run variability, calibration, and balance uncertainty. A null is informative only if its interval excludes the consequential benefit. PDVB remains present during recovery, so improved retention demonstrates the function of the complete formulation, not the number of surviving cobalt sites or permanent irreversibility.

A positive contrast after late addition excludes an explanation confined entirely to the original conditioning, but it leaves effects during restart, the water transient, and ongoing liquid or chemical changes open. Any parent checks used to reject a whole-bed reaction-water source explanation must bracket the measured formulation-dependent water histories during both challenge and recovery, including their uncertainty. Plateau-only bracketing cannot exclude that explanation for post-recovery retention: generated water has a larger relative contribution during recovery. If the relevant histories cannot be bounded, retain the source explanation while reporting the formulation contrast. A null, after successful reproduction and qualification, limits the supported benefit of late addition in this specific state; it does not establish that startup protection caused the published result.

\subsection{Experiment 3: predict one useful disturbance response}

Pursue a predictive branch only if the preceding comparison or reversible water response supplies a resolvable functional question. Zero-length-column work shows how independent storage information and flow-dependent transients can constrain water exchange in silica gel, but its passive, low-temperature coefficients cannot be transferred to a reacting, liquid-bearing cobalt bed.\cite{brandani2022}

First qualify water-specific blanks at the actual temperature, pressure, and composition, including heated transfer lines, pressure reduction, and detector memory. An inert tracer establishes gas delay but cannot correct water-specific storage. On representative conditioned parent and PDVB packings, apply small upward and downward humidity steps at two flows with unchanged inlet H\textsubscript{2}/CO partial pressures, a reduced-amplitude repeat, and baseline returns. Integrate absolute water rates; do not normalize away the inventory information. Measure product rates fast enough to bound the changing chemical water source. A total-mass uptake decrease must also be corrected for dilution: for additive/catalyst mass ratio $a$, noninteracting uptake is $(u_{\mathrm{cat}}+a u_{\mathrm{add}})/(1+a)$ even with unchanged catalyst loading.

Fit only the simplest response representation consistent with both step directions, flows, source balance, and independent packings. In a reacting compartment, linearization of Eq.~\ref{eq:water-pool} with $q=q(c)$ gives a relaxation rate $(G-\partial q/\partial c)/B$. This decay expression assumes $B>0$ and a stable stationary state with $G>\partial q/\partial c$, the derivative evaluated at that state; a measured pole is not automatically $G/B$. If source feedback, blank uncertainty, or changing liquid state prevents a meaningful coefficient, stop the conductance claim. Empirical input--output prediction can still be tested, but its coefficients must not be relabeled local transport properties.

The concrete validation is a partial-duration water pulse withheld from fitting. Preselect a fixed decision horizon from pulse onset, representing the available production window, and an acceptable C\textsubscript{5+} carbon-output deficit relative to undisturbed operation. This endpoint informs whether to tolerate a humidity excursion of the proposed duration within that output budget. Calibrate each formulation's nearby reversible product response independently and account for product-detector and liquid-collection delays. Select a pulse duration for which surviving models predict consequentially different output deficits over that fixed horizon, and freeze predictions before collecting it. Distinct water traces alone do not qualify: different time constants with equal steady gains can yield the same deficit integrated through complete recovery.

Predict both the water trace and the selected functional endpoint, with an endpoint prediction interval narrower than the predeclared consequential contrast. Compare with a simple baseline using inlet humidity, the calibrated steady product response, and independently measured apparatus delays; this uses the same calibration and validation data. Validate on independent charges without refitting parameters to the withheld pulse. Passing this test supports a bounded operating-history rule over the calibrated state and amplitude range. If transient characterization does not improve on the baseline, claim only the supported empirical rule. Neither outcome validates a damage law at much higher water exposure.

\subsection{Decision and scope of the original contribution}

The program has two credible contributions: a controlled test of late low-dose PDVB at a defined catalyst state and disturbance, and, conditionally, a prospective rule connecting a measured disturbance to useful catalytic output within the calibrated range. Their importance is deciding whether a simple formulation change pays for its initial penalty under the tested conditions and when faster water clearance is an inadequate selection criterion. A mechanism paper would require an additional selective constraint on the producing region, shown to remain valid when the additive changes liquid and catalyst state. Outlet transients alone cannot provide that guarantee.

Close the affected branch if the source effect cannot be reproduced, handling destroys interpretability, the challenge has no measurable dynamic range, or the predicted functional difference is below measurement resolution. If fractional protection is real but cumulative output is worse, report that tradeoff and do not claim a practical advance. A useful formulation result can survive failure to identify the stationary mechanism. Conversely, a well-fitted water curve without a consequential functional prediction does not justify an expanded materials or isotope campaign. Commercial claims require a subsequent complete-cycle comparison with an optimized catalyst, including regeneration, additive inventory, bed volume, and replacement time.
